
Can a language model's architecture family be identified from its adversarial failures alone? Current robustness assessments treat models as individuals, aggregating vulnerability into a scalar and discarding the multivariate structure across attack types that architecture constrains. This paper introduces the $\Delta$*-vector framework, which represents each model as a profile of normalized per-attack-category vulnerability scores and tests whether architecture family (encoder-only, decoder-only, encoder-decoder) explains this profile. Applied to nine scale-matched transformer models across six adversarial attack categories (AdvGLUE and ANLI-R3), permutation MANOVA yields $\eta^{2}=0.293$ --- a large effect (Cohen's $f^{2}\approx 0.41)$ that exceeds the pre-specified $\eta^{2}=0.15$ threshold in five of six attack categories (83\%). Statistical significance at $\alpha =0.05$ is not achieved (p=0.147, 1,000 permutations), consistent with the approximately 40\% statistical power available at N=9; $N\geq 15$ ($\geq 5$ per family) is required for 80\% power. Leave-one-model-out (LOMO) classification achieves accuracy=0.333 --- at three-class chance --- a result attributable to geometric degeneracy at N=3 models per family rather than absence of family structure. A validated seven-module pipeline (1,788 lines, Python) is provided. The existence signal is confirmed at effect-size level; overall statistical significance awaits a confirmatory study with an expanded model pool.

